\documentclass[11pt]{article}

\usepackage[a4paper,margin=1in]{geometry}
\usepackage{amsmath,amssymb,amsthm,mathtools}
\usepackage[hidelinks]{hyperref}
\hypersetup{
  pdfauthor={Bangyen Pham},
  pdftitle={Light Integer Multiples at Gaussian Roots: Congruences, Level Sets and Pell Points},
  pdfsubject={Integer polynomial multiples with prescribed Gaussian-integer roots and small coefficients},
  pdfkeywords={polynomial multiples, Gaussian integers, level sets, p-adic valuations, Pell equation},
  pdflang={en-US}
}

\newtheorem{theorem}{Theorem}[section]
\newtheorem{lemma}[theorem]{Lemma}
\newtheorem{corollary}[theorem]{Corollary}
\newtheorem{proposition}[theorem]{Proposition}
\newcommand{\R}{\mathbb R}
\newcommand{\N}{\mathbb N}

\title{Light Integer Multiples at Gaussian Roots: Congruences, Level Sets and Pell Points}
\author{\shortstack{Bangyen Pham\\
\small Independent Researcher\\
\small\href{mailto:bangyenp@gmail.com}{\texttt{bangyenp@gmail.com}}\\
\small\href{https://orcid.org/0009-0006-9928-2088}{ORCID: 0009-0006-9928-2088}}}
\date{September 2026}

\begin{document}
\maketitle

\begin{abstract}
Let $P=\prod_j(x-\alpha_j)(x-\bar\alpha_j)$ for distinct Gaussian integers
$\alpha_1,\ldots,\alpha_K$ above the real axis, of least modulus
$\rho_{\min}$.  We study integer multiples of $P$ whose coefficients are
small.  The lowest $m$ coefficients of the multiples are cut out by one
congruence modulo $P(0)^m$, which gives lower bounds for the lowest
coefficient and lines of root sets along which $P'(0)$ is constant.  If
$\rho_{\min}>2$ and the second largest nonleading coefficient is below
$\rho_{\min}/2$, all the $\alpha_j$ lie in one level set $H(y)=-v$ of an
integer polynomial $H$ with small nonleading coefficients.  Such a level set
lies near one circle, so $K<6\rho_{\min}$, and three Gaussian points
above the axis share one at arbitrarily large $\rho_{\min}$, on the
solutions of a Pell equation.  At each prime the lowest term of $H$ controls
the roots, so at most $\operatorname{ord}_0H$ of them share a norm; a unit
relating two of them is a symmetry of $H$, and four points of one class
force $\deg H\ge16$.
\end{abstract}

\noindent\textbf{Keywords.} Polynomial multiples; Gaussian integers; level
sets; $p$-adic valuations; Pell equation.

\noindent\textbf{2020 Mathematics Subject Classification.} Primary 11C08; Secondary 11R04, 11D09.

\section{Introduction}\label{sec:setting}

A real multiple of a polynomial with $K$ prescribed roots of modulus at
least $R$ in a sector can have logarithmic coefficient mass as small as
order $K\log R$: $x^N-\rho^N$ has only two nonzero coefficients.  Whether
integer multiples at Gaussian-integer roots must pay order $K^2\log R$ is
open~\cite{coefficient-mass-sectors}, and every integer imitation of
$x^N-\rho^N$ is expensive there.  This paper studies the multiples that
come closest: integer multiples whose coefficients other than the lowest
are small.

\paragraph{Setting.}  Throughout, $\alpha_1,\ldots,\alpha_K\in\mathbb Z[i]$
are distinct, with $\alpha_j=a_j+ic_j$ and $c_j\ge1$; put
$\rho_j=|\alpha_j|$, $\rho_{\min}=\min_j\rho_j$,
$\rho_{\max}=\max_j\rho_j$ and
\[
  P=\prod_{j=1}^K(x-\alpha_j)(x-\bar\alpha_j)
   =\prod_{j=1}^K\bigl(x^2-2a_jx+a_j^2+c_j^2\bigr)\in\mathbb Z[x],
\]
a monic polynomial with $2K$ distinct roots and $P(0)\ne0$.  For
$F=\sum_{i=0}^Df_ix^i$ a \emph{split} is an $m\in\{1,\ldots,D\}$, with
$F=F_{<m}+x^mG_m$, $F_{<m}=\sum_{i<m}f_ix^i$; it is \emph{clean at} $\alpha$
if $G_m(\alpha)=0$.  We use three results of~\cite{coefficient-mass-sectors}.
If $F(\alpha)=0$, $\rho=|\alpha|>2$ and the split $m$ is not clean at
$\alpha$, some $i<m$ has $|f_i|\ge(\rho/2)^{m-i}$
\cite[Lemma~3.1]{coefficient-mass-sectors}.  Every nonzero integer multiple
of $P$ is a sum $\sum_sx^{e_s}B_s$ of \emph{blocks} with disjoint supports,
each an integer multiple of $P$ with no split clean at every $\alpha_j$
\cite[Lemma~3.2]{coefficient-mass-sectors}.  At odd, pairwise coprime norms
$n_j=\rho_j^2$ with $\gcd(a_j,c_j)=1$, a gap $m$ after the lowest nonzero
coefficient $f_e$ forces $|f_e|\ge P(0)^m$
\cite[Proposition~5.3]{coefficient-mass-sectors}; its item~2 and the
lattice of \cite[Proposition~5.4]{coefficient-mass-sectors} are recalled
where they are used.

\paragraph{Results.}  The lowest $m$ coefficients of the multiples are cut
out by one congruence modulo $P(0)^m$, at a root of $P$ found by Newton's
method from $0$ (Proposition~\ref{prop:gausscong}); along lines of root sets
with a constant $P'(0)$ (Proposition~\ref{prop:gaussline}) this bounds what
the lowest coefficient can gain from small coefficients after it.  A
multiple is \emph{light} if its coefficients other than the lowest are below
$\rho_{\min}/2$; light multiples have degree at least $2K$
(Proposition~\ref{prop:gausslight}).  If $\rho_{\min}>2$ and
$b_2(F)<\rho_{\min}/2$, then $F=x^e(v+H)$ with all the $\alpha_j$ in the
level set $H(y)=-v$ (Proposition~\ref{prop:gaussheavy}).  The $\alpha_j$ then
lie near one circle, so $K<6\rho_{\min}$
(Proposition~\ref{prop:gausscircle}), and three Gaussian points above the
axis share such a level set at arbitrarily large $\rho_{\min}$.  At each
prime the lowest term $h_\ell y^\ell$ of $H$ controls the roots, and at most
$\ell$ of them share a norm (Proposition~\ref{prop:gaussprimes}); a unit
relating two of them is a symmetry of $H$ (Proposition~\ref{prop:gausssym}),
and four points of one class need $\deg H\ge16$
(Proposition~\ref{prop:gaussfour}).  No proof relies on computation.

\paragraph{Open.}  Whether the number of Gaussian points above the axis in
one level set of a polynomial with small nonleading coefficients is bounded
independently of $\rho_{\min}$, and whether one norm carries at most two of
them.

\paragraph{Conventions.}  All logarithms are natural.  For a polynomial
$F$, $b_k(F)$ is its $k$-th largest nonleading coefficient magnitude and
$\Lambda(F)=\sum_j\log^+|f_j|$ its logarithmic coefficient mass.
References of the form \cite[Lemma~3.1]{coefficient-mass-sectors} are to
the companion paper.

\section{One congruence}\label{sec:congruence}

The lowest $m$ coefficients of the multiples are cut out by one congruence,
modulo $P(0)^m$, at a root of $P$ found by Newton's method from $0$; this
holds for every monic $P$ with $\gcd(P(0),P'(0))=1$, and so recovers
\cite[Proposition~5.3]{coefficient-mass-sectors}, item~2.

\begin{proposition}[One congruence]\label{prop:gausscong}
Let $P\in\mathbb Z[x]$ be monic with $D=P(0)\ne0$, and let $m\ge1$.
\begin{enumerate}
\item If $\gcd(D,P'(0))=1$, there is $t\in\mathbb Z$ with $D\mid t$ and
$P(t)\equiv0\pmod{D^m}$, and for $(f_0,\ldots,f_{m-1})\in\mathbb Z^m$ some
$F\in\mathbb Z[x]$ divisible by $P$ has $F\equiv\sum_{i<m}f_ix^i\pmod{x^m}$
if and only if
\[
  \sum_{i<m}f_it^i\ \equiv\ 0\pmod{D^m}.
\]
In particular $d_m=D^m$ in the notation of \cite[Proposition~5.3]{coefficient-mass-sectors}.
In the setting of Section~\ref{sec:setting} with odd pairwise coprime norms
$n_j$ and $\gcd(a_j,c_j)=1$, $\gcd(D,P'(0))=1$.
\item Some $F\in\mathbb Z[x]$ divisible by $P$ has $F\equiv f_0+f_1x\pmod{x^2}$
if and only if $D\mid f_0$ and $f_1\equiv P'(0)f_0/D\pmod D$.
\item In the setting of Section~\ref{sec:setting}, if
$|P'(0)|<\rho_{\min}/2$, then $F=P$ has $|p_1|<\rho_{\min}/2$ and
$\log|p_0|=\log P(0)$, so a constant $C$ with
$\log|f_e|\ge2\log P(0)-2C\log\rho_{\max}$ for every integer multiple with
$|f_{e+1}|<\rho_{\min}/2$ has $C\ge\log P(0)/(2\log\rho_{\max})$.  At the
four points $-26+35i$, $-1+60i$, $29+62i$, $47+62i$, of norms $1901$,
$3601=13\cdot277$, $4685=5\cdot937$, $6053$, which are odd and pairwise
coprime, with $\gcd(a_j,c_j)=1$ and $43<\rho_j<78$, $P'(0)=6$, and so
$C>3.77$.
\end{enumerate}
\end{proposition}

\begin{proof}
1.  Put $t_1=0$, and given $t_k$ with $D\mid t_k$ and $D^k\mid P(t_k)$, note
$P'(t_k)\equiv P'(0)\pmod D$, so there is $w\in\mathbb Z$ with
$wP'(t_k)\equiv1\pmod{D^m}$; put $h=-wP(t_k)$ and $t_{k+1}=t_k+h$.  Then
$D\mid t_{k+1}$, and $P(t_k+h)=P(t_k)+P'(t_k)h+h^2Q$ with $Q\in\mathbb Z$,
so $P(t_{k+1})\equiv P(t_k)(1-wP'(t_k))\equiv0\pmod{D^{k+1}}$, as
$D^k\mid P(t_k)$ and $D^k\mid h$.  Take $t=t_m$.  If $F=PS$ with
$S\in\mathbb Z[x]$ and $F=\sum_{i<m}f_ix^i+x^mG$, then
$\sum_{i<m}f_it^i=P(t)S(t)-t^mG(t)\equiv0\pmod{D^m}$, as $D^m\mid t^m$.  So
the lattice $L_m$ of the proof of \cite[Proposition~5.4]{coefficient-mass-sectors}, item~3,
lies in the kernel of $f\mapsto\sum_{i<m}f_it^i\bmod D^m$, a homomorphism
$\mathbb Z^m\to\mathbb Z/D^m$ that is onto, as $f=(r,0,\ldots,0)$ maps to
$r$.  Both have index $D^m$, so they are equal, and $(f,0,\ldots,0)$ lies in
$L_m$ exactly when $D^m\mid f$.  In the setting of Section~\ref{sec:setting},
$P=\prod_j(x^2-2a_jx+n_j)$ gives $P'(0)=-\sum_j2a_jD/n_j$, and a prime $p$
dividing $D$ divides exactly one $n_j$, is odd and divides neither $a_j$
(otherwise $p\mid\gcd(a_j,c_j)$) nor $D/n_j$, so
$P'(0)\equiv-2a_jD/n_j\not\equiv0\pmod p$.

2.  In the notation of the proof of \cite[Proposition~5.4]{coefficient-mass-sectors}, item~3,
$L_2$ is spanned by $\mathbf p_0=(D,P'(0))$ and $\mathbf p_1=(0,D)$, so its
vectors are $(sD,sP'(0)+s'D)$ with $s,s'\in\mathbb Z$.

3.  $P$ is divisible by $P$, with $p_0=P(0)$ and $p_1=P'(0)$; the inequality
at $F=P$, $e=0$, reads $\log P(0)\ge2\log P(0)-2C\log\rho_{\max}$.  For the
example, $P'(0)=-2\sum_ja_jD/n_j=6<\rho_{\min}/2$, as $\rho_{\min}^2=1901$,
and $\log D/(2\log\rho_{\max})>3.77$ with $D=194126805235805$ and
$\rho_{\max}^2=6053$.
\end{proof}

So at $m=2$ the divisibility $P(0)\mid f_e$ is, at such root sets, all that
holds: a bound $\log|f_e|\ge m\log P(0)-Cm\log\rho_{\max}$ for all
integer multiples with $|f_i|<\rho_{\min}/2$ for $e<i<e+m$ needs, at $m=2$,
$C\ge\log P(0)/(2\log\rho_{\max})$, which is nearly $K$ for points in one
annulus, whenever the multiple $P$ itself has $|P'(0)|<\rho_{\min}/2$, that
is, whenever $\bigl|\sum_ja_j/n_j\bigr|<\rho_{\min}/(4P(0))$.  For $K\le3$
this happens at infinitely many root sets in one annulus, along lines on
which $P'(0)$ does not move (Proposition~\ref{prop:gaussline}); whether it
happens for every $K$ is open.  The multiples of
Proposition~\ref{prop:gaussheavy}, item~2, need the step only at
$m=\deg H\ge2K$, where the congruence of item~1 involves $m-1\ge2K-1$ small
coefficients.

\begin{proposition}[A constant derivative along lines]\label{prop:gaussline}
\begin{enumerate}
\item Let $\lambda_j,\mu_j\in\mathbb Z[i]$, $\lambda_j\ne0$, be such that the
$2K$ numbers $w_j=-\mu_j/\lambda_j$ and $\bar w_j$ are distinct, and for real
$s$ let $P_s=\prod_j(x-\alpha_j)(x-\bar\alpha_j)$ with
$\alpha_j=\lambda_js+\mu_j$.  Put $\Lambda=\prod_j|\lambda_j|^2$ and
$Q(s)=\prod_j(s-w_j)(s-\bar w_j)$.  Then $P_s(0)=\Lambda Q(s)$, and $P_s'(0)$
is a polynomial in $s$ with integer coefficients and degree less than $2K$,
equal to a constant $-N$, $N\in\mathbb R$, if and only if $N\lambda_j=\Lambda Q'(w_j)$ for
every $j$.
\item In the setting of Section~\ref{sec:setting}, at $\alpha_1=-(s+3)+(s+2)i$
and $\alpha_2=(s+1)+(s+2)i$,
$P'(0)=4$ for every $s$, and for every integer $s\ge-1$ the norms are odd and
coprime and $\gcd(a_j,c_j)=1$.
\item At $\alpha_1=-(s+3)+(s+2)i$, $\alpha_2=-1+(s+1)i$ and
$\alpha_3=(s-1)+si$, $P'(0)=90$ for every $s$, and for every odd integer
$s\ge1$ with $s\equiv0$ or $3\pmod5$ the norms are odd and pairwise coprime and
$\gcd(a_j,c_j)=1$.
\end{enumerate}
\end{proposition}

\begin{proof}
1.  As polynomials in $s$, $\bar\alpha_j=\bar\lambda_js+\bar\mu_j$ and
$\alpha_j\bar\alpha_j=|\lambda_j|^2(s-w_j)(s-\bar w_j)$, so
$P_s(0)=\prod_j\alpha_j\bar\alpha_j=\Lambda Q(s)$, and
$\Phi(s)=P_s'(0)=-\sum_j(\alpha_j+\bar\alpha_j)\prod_{l\ne j}\alpha_l\bar\alpha_l$
has integer coefficients and degree at most $2K-1$.  As rational functions,
\[
  \frac{\Phi(s)}{\Lambda Q(s)}=-\sum_j\Bigl(\frac1{\lambda_j(s-w_j)}
  +\frac1{\bar\lambda_j(s-\bar w_j)}\Bigr),
\]
and the $2K$ poles are simple, so the residues at $w_j$ give
$\Phi(w_j)=-\Lambda Q'(w_j)/\lambda_j$, and $\Phi(\bar w_j)$ is its conjugate,
$\Phi$ and $Q$ having real coefficients.  A polynomial of degree less than
$2K$ is the constant $-N$ if and only if it takes the value $-N$ at the $2K$
distinct points $w_j,\bar w_j$, that is, for real $N$, if and only if
$\Lambda Q'(w_j)=N\lambda_j$ for every $j$.

2.  Here $a_1=-(s+3)$, $a_2=s+1$, $n_1=2s^2+10s+13$ and $n_2=2s^2+6s+5$, and
expanding, $a_1n_2+a_2n_1=-2$, so $P'(0)=-2(a_1n_2+a_2n_1)=4$.  Both norms
are odd; $n_1-n_2=4(s+2)$ and $n_2=2(s+2)^2-2(s+2)+1\equiv1\pmod{s+2}$, so
$\gcd(n_1,n_2)=1$; and $\gcd(s+3,s+2)=\gcd(s+1,s+2)=1$.

3.  Here $n_1=2s^2+10s+13$, $n_2=s^2+2s+2$, $n_3=2s^2-2s+1$, and expanding,
$P'(0)=-2(a_1n_2n_3+a_2n_1n_3+a_3n_1n_2)=90$.  The real and imaginary parts
of each $\alpha_j$ are coprime, being $1$ or consecutive integers, so no
prime $p\equiv3\pmod4$ divides a norm; $n_1$ and $n_3$ are odd, and $n_2$ is
odd as $s$ is.  Next
\[
  n_1-n_3=12(s+1),\qquad n_1-2n_2=3(2s+3),\qquad 2n_2-n_3=3(2s+1),
\]
while $n_3=2(s+1)^2-6(s+1)+5$ and $4n_2=(2s+1)(2s+3)+5$.  So each
$\gcd(n_j,n_l)$ divides $5$, and $5$ divides one only if $5$ divides $s+1$,
$2s+3$ or $2s+1$, that is, $s\equiv4$, $1$ or $2\pmod5$.
\end{proof}

Along items~2 and~3, $|P'(0)|<\rho_{\min}/2$ once $\rho_{\min}>8$, respectively
$\rho_{\min}>180$, and $\rho_{\max}/\rho_{\min}$ tends to $1$, respectively
$\sqrt2$, so for large $s$ the points lie in one annulus $[R,2R]$, and
$\log P(0)/(2\log\rho_{\max})$ tends to $K=2$, respectively $3$; so a
constant $C$ as in Proposition~\ref{prop:gausscong}, item~3, valid along
them is at least $K$.  More generally, lines as in item~1 with a constant
$P'(0)=-N$ that meet the hypotheses of \cite[Proposition~5.3]{coefficient-mass-sectors},
item~2, at infinitely many integers $s$ (with $\bar\alpha_j$ in place of
$\alpha_j$ where $\operatorname{Im}\alpha_j<0$) force $C\ge K$, since
$\rho_j=|\lambda_j||s-w_j|$ and $P(0)=\Lambda Q(s)$, in one annulus when the
$|\lambda_j|$ are within a factor less than $2$ of each other.  In item~3,
$w_1=-\frac{5+i}2$, $w_2=-1-i$, $w_3=\frac{1-i}2$, and $\lambda_j$ is the
denominator of $w_j$ up to a unit, $\Lambda=4$ and $N=-90$.

\section{Light multiples and level sets}\label{sec:light}

The step at $m\ge2K$ contains a question about level sets.  Call a nonzero
integer multiple $F$ of $P$ \emph{light} if every coefficient of $F$ other than
its lowest nonzero one has modulus less than $\rho_{\min}/2$.

\begin{proposition}[Light multiples]\label{prop:gausslight}
In the setting of Section~\ref{sec:setting}, assume $\rho_{\min}\ge2$, and
let $F$ be a light multiple of $P$, of degree $e+d$ with lowest nonzero
coefficient $f_e$.  Then $d\ge2K$ and $|f_e|<\rho_{\min}^{d+1}$.  So a constant
$C$ with $\log|f_e|\ge m\log P(0)-Cm\log\rho_{\max}$ at $m=d+1$ satisfies
\[
  C\log\rho_{\max}\ >\ \log P(0)-\log\rho_{\min}\ \ge\ (2K-1)\log\rho_{\min}.
\]
\end{proposition}

\begin{proof}
As $P(0)\ne0$, $P$ divides $F/x^e$, a nonzero polynomial of degree $d$, so
$d\ge2K$.  Take $j$ with $\rho_j=\rho_{\min}=\rho$.  From
$F(\alpha_j)=0$, $f_e=-\sum_{i=1}^df_{e+i}\alpha_j^i$, so
\[
  |f_e|\ <\ \frac\rho2\sum_{i=1}^d\rho^i\ <\ \frac\rho2\cdot\frac{\rho^{d+1}}{\rho-1}
  \ \le\ \rho^{d+1},
\]
as $\rho\ge2$.  The coefficients $f_i$, $e<i<e+d+1$, have modulus less than
$\rho/2$, so the bound at $m=d+1$ applies, and
$(d+1)(\log P(0)-C\log\rho_{\max})\le\log|f_e|<(d+1)\log\rho$.  Finally
$P(0)=\prod_j\rho_j^2\ge\rho^{2K}$.
\end{proof}

A light multiple is a multiple of Proposition~\ref{prop:gaussheavy}, item~2,
whose leading coefficient $h_d$ is also below $\rho_{\min}/2$; the three-point
level sets after that proposition, with $h_d=1$, are light, at a shared norm.
So an affirmative answer to the step at every $m\ge2K$, with a constant $C$,
at odd pairwise coprime norms with $\gcd(a_j,c_j)=1$, bounds the number of
points of such level sets:
$K<\frac12(C\log\rho_{\max}/\log\rho_{\min}+1)$, which in one annulus
$[R,2R]$ is $\frac12(C+1)+O(C/\log R)$.  Whether that number is bounded is
open (see the end of this section), here for light multiples at coprime
norms, and no proof of the step at $m\ge2K$ can avoid it.

At Gaussian roots a small second row confines an integer multiple to one
level set of a polynomial with small nonleading coefficients, whereas over
the reals $x^N-\rho^N$ has $b_2=0$.

\begin{proposition}[One heavy coefficient]\label{prop:gaussheavy}
In the setting of Section~\ref{sec:setting}, assume $\rho_{\min}>2$, and let
$F\in\mathbb Z[x]$ be nonzero and divisible by $P$, with lowest nonzero
coefficient $v=f_e$.  Call a coefficient $f_i$ of $F$ \emph{heavy} if
$|f_i|\ge\rho_{\min}/2$.
\begin{enumerate}
\item For any decomposition of $F$ as in \cite[Lemma~3.2]{coefficient-mass-sectors}, the
lowest coefficient of every block is heavy and nonleading, so there are at most
as many blocks as heavy nonleading coefficients of $F$.
\item If $b_2(F)<\rho_{\min}/2$, then $F=x^e(v+H)$ with $H\in\mathbb Z[x]$,
$H(0)=0$ and every nonleading coefficient of $H$ of modulus less than
$\rho_{\min}/2$; moreover $H(\alpha_j)=-v$ for every $j$, and
$(\rho_{\min}/2)^{\deg H}\le|v|$.
\end{enumerate}
\end{proposition}

\begin{proof}
1.  A block $B$ is $PQ$ with $Q\in\mathbb Z[x]$, and $B(0)\ne0$, so
$|B(0)|\ge P(0)\ge\rho_{\min}^2>\rho_{\min}/2$.  Its degree is at least $2K\ge2$,
so its lowest coefficient is not the leading coefficient of $F$, and distinct
blocks have distinct lowest positions.

2.  At most one nonleading coefficient of $F$ is heavy, so by item~1 a
decomposition of $F$ as in \cite[Lemma~3.2]{coefficient-mass-sectors} has exactly one block,
$F=x^eB$, and $B(0)=v$ is the heavy one; every other
nonleading coefficient of $F$ has modulus less than $\rho_{\min}/2$.  Write
$B=v+H$; since $B(\alpha_j)=0$, $H(\alpha_j)=-v$.  Let
$d=\deg H=\deg B$.  The split $d$ of the $P$-irreducible $B$ is unclean at
some $\alpha_j$, and $\rho_j\ge\rho_{\min}>2$, so by
\cite[Lemma~3.1]{coefficient-mass-sectors} some $i<d$ has
$|b_i|\ge(\rho_j/2)^{d-i}\ge\rho_{\min}/2$.  This $b_i$ is a heavy
nonleading coefficient of $F$, so $i=0$ and $|v|\ge(\rho_{\min}/2)^d$.
\end{proof}

So, when $\rho_{\min}>2$, $b_2(F)\ge\rho_{\min}/2$ unless every $\alpha_j$
lies in one level set $\{y:H(y)=-v\}$ of such an $H$.  Three Gaussian points
above the axis can share such a level set, at arbitrarily large
$\rho_{\min}$.  Let $a^2-3c^2=1$ with $a>0$ and $c\ge4$; these solutions are
$a+c\sqrt3=(2+\sqrt3)^n$, $n\ge2$, so $c=4,15,56,\ldots$.  Then
\[
  y^3-y-(8c^3+2c)i=(y+2ci)\bigl(y^2-2ciy-(4c^2+1)\bigr)
  =(y+2ci)(y-a-ci)(y+a-ci),
\]
so $y^3-y$ is $(8c^3+2c)i$ at $\pm a+ci$ and $-(8c^3+2c)i$ at $2ci$, and
$(y^3-y)^2$, with nonleading coefficients $-2$ and $1$, takes the value
$-(8c^3+2c)^2$ at the three points $\pm a+ci$, $2ci$, which have
$\rho_{\min}=2c>4$.  At $c=4$ these are $\pm7+4i$, $8i$ and the value is
$-270400$.  In the other direction, the $\alpha_j$ lie near one circle, so
there are fewer than $6\rho_{\min}$ of them.

\begin{proposition}[Near one circle]\label{prop:gausscircle}
In the setting of Proposition~\ref{prop:gaussheavy}, item~2, let
$d=\deg H$ and $Q=(3\rho_{\min}-2)/(\rho_{\min}-2)$.  Then
$\rho_{\min}\le\rho_j<\rho_{\min}Q^{1/d}$ for every $j$, and
$K<6\rho_{\min}$.  Consequently
$\Lambda(F)\ge2K\log(\rho_{\min}/2)>2K\log(K/12)$.
\end{proposition}

\begin{proof}
Write $\rho=\rho_{\min}$, $H=\sum_{i=1}^dh_iy^i$ and
$B=\max_{0<i<d}|h_i|<\rho/2$, and put $\kappa=\rho/(2(\rho-1))$, so
$\kappa<1$ as $\rho>2$, and $(1+\kappa)/(1-\kappa)=Q$.  For each $j$, with
$r=\rho_j\ge\rho$, $H(\alpha_j)=-v$ gives
\[
  |h_d\alpha_j^d+v|=\Bigl|\sum_{0<i<d}h_i\alpha_j^i\Bigr|
  \le B\,\frac{r^d-r}{r-1}<\frac{\rho}{2(\rho-1)}\,r^d\le\kappa|h_d|r^d,
\]
so $(1-\kappa)|h_d|\rho_j^d<|v|<(1+\kappa)|h_d|\rho_j^d$.  The upper bound at a
$j$ with $\rho_j=\rho$ and the lower bound at any $j$ give
$(\rho_j/\rho)^d<Q$.

Since $\rho^2$ is an integer above $4$, $\rho\ge\sqrt5$, so
$Q=3+4/(\rho-2)\le11+4\sqrt5<20$ and $\log Q<3<\frac32\rho$.  Suppose
$K\ge6\rho$.  The $2K$ distinct roots $\alpha_j,\bar\alpha_j$ of $v+H$ give
$d\ge2K\ge12\rho$, so $t=(\log Q)/d$ satisfies $t<\frac18$ and
$t<\frac3{12\rho}$.  Put $\rho'=\rho Q^{1/d}=\rho e^t$.  The unit squares
centred at the $\alpha_j$ are disjoint and lie in the upper half of the
annulus $\rho-1/\sqrt2\le|w|\le\rho'+1/\sqrt2$, so
\[
  K\ \le\ \frac\pi2\Bigl(\rho'^2-\rho^2+\sqrt2(\rho+\rho')\Bigr)
  =\frac\pi2\rho\Bigl(\rho(e^{2t}-1)+\sqrt2(1+e^t)\Bigr).
\]
Here $e^{2t}-1\le2te^{2t}<2.6\,t$ and $e^t<1.14$, so
$\rho(e^{2t}-1)<2.6\cdot\frac3{12}<0.66$ and
$K<\frac\pi2\rho(0.66+3.03)<5.8\rho$, a contradiction.  Finally
$\Lambda(F)\ge\log|v|\ge d\log(\rho/2)\ge2K\log(\rho/2)$ by
Proposition~\ref{prop:gaussheavy}, and $\rho/2>K/12$.
\end{proof}

\section{The arithmetic of a level set}\label{sec:levelprimes}

The count uses only the annulus, whose width is of order $\rho_{\min}/d$,
and not the integrality of $H$; when $d\ge8\rho_{\min}^2$ it gives
$\rho_j^2<\rho_{\min}^2+1$, so all the $\alpha_j$ share one norm.  The
integrality of $H$ acts at each prime, through the lowest term of $H$.  For
a prime $p$ let $\operatorname{ord}_p$ be the valuation of $\mathbb C_p$ with
$\operatorname{ord}_p(p)=1$, and fix an embedding of $\mathbb Q(i)$ in
$\mathbb C_p$.

\begin{proposition}[Primes of a level set]\label{prop:gaussprimes}
In the setting of Proposition~\ref{prop:gaussheavy}, item~2, write
$H=\sum_{i=1}^dh_iy^i$ with $h_d\ne0$, let $\ell$ be the least $i\ge1$ with
$h_i\ne0$, put $n_j=\rho_j^2$, and for a prime $p$ put
$t_p=\operatorname{ord}_p(h_\ell)$ and
$s_p=(\operatorname{ord}_p(v)-t_p)/\ell$.
\begin{enumerate}
\item For every prime $p$, at most $\ell$ roots of $v+H$ in $\mathbb C_p$,
counted with multiplicity, have valuation greater than $t_p$, and each of
them has valuation $s_p$.
\item Let $\ell<d$.  For every prime $p$, at most $\ell$ of the $\alpha_j$ have
$\operatorname{ord}_p(n_j)>2t_p$, and at most $\ell/2$ if $p=2$ or
$p\equiv3\pmod4$; each of them has $\operatorname{ord}_p(n_j)=2s_p$ or
$s_p\le\operatorname{ord}_p(n_j)\le s_p+t_p$.  Every $n_j$ has such a prime,
so at most $\ell$ of the $\alpha_j$ share one norm, and $K\le\ell N$ if the
$n_j$ take $N$ values.
\item If all the $\alpha_j$ share one norm, in particular if
$d\ge8\rho_{\min}^2$, then $K\le\ell$ when $\ell<d$, and $K\le2$ when $\ell=d$.
\end{enumerate}
\end{proposition}

\begin{proof}
1.  Let $V=\operatorname{ord}_p(v)$.  The Newton polygon of $v+H$ over
$\mathbb Q_p$ is the lower convex hull of $(0,V)$ and the points
$(i,\operatorname{ord}_p(h_i))$ with $h_i\ne0$; none has $0<i<\ell$, and all
have ordinate at least $0$, the coefficients being integers.  For each edge
of slope $-\sigma$ and horizontal length $\lambda$, exactly $\lambda$ roots
in $\mathbb C_p$, counted with multiplicity, have valuation $\sigma$
\cite[Chapter~II, Proposition~6.3]{neukirch}.  Let the first edge run from
$(0,V)$ to $(i_1,y_1)$ with slope $-\sigma_1$, so $i_1\ge\ell$ and $y_1\ge0$.
If $i_1>\ell$, the point $(\ell,t_p)$ lies on or above this edge, so
$t_p\ge y_1+\sigma_1(i_1-\ell)\ge\sigma_1$, and by convexity no root has
valuation above $\sigma_1\le t_p$.  If $i_1=\ell$, the first edge carries
$\ell$ roots of valuation $(V-t_p)/\ell=s_p$.  Every later edge of negative
slope starts at an abscissa at least $\ell$, where the polygon, decreasing
from $(\ell,t_p)$, has ordinate at most $t_p$; it has horizontal length at
least $1$ and ends at ordinate at least $0$, so its roots have valuation at
most $t_p$.

2.  Fix $p$ and $j$ with $e=\operatorname{ord}_p(n_j)>2t_p$.  The Gaussian
integers $\alpha_j$ and $\bar\alpha_j$ have nonnegative valuations with sum
$\operatorname{ord}_p(\alpha_j\bar\alpha_j)=e$, so one of them exceeds
$t_p$.  If $p=2$ or $p\equiv3\pmod4$, $-1$ is not a square in $\mathbb Q_p$,
the nontrivial automorphism of the quadratic extension $\mathbb Q_p(i)$
restricts to complex conjugation on $\mathbb Q(i)$, and the valuation, being
unique on $\mathbb Q_p(i)$, is invariant under it; so
$\operatorname{ord}_p(\alpha_j)=\operatorname{ord}_p(\bar\alpha_j)=e/2>t_p$.
The $2K$ numbers $\alpha_j,\bar\alpha_j$ are distinct roots of $v+H$, so
item~1 gives the two counts.  By item~1 again each of
$\alpha_j,\bar\alpha_j$ of valuation above $t_p$ has valuation $s_p$: if
both do, $e=2s_p$, and otherwise one has valuation $s_p$ and the other at
most $t_p$, so $s_p\le e\le s_p+t_p$.  Since $\ell<d$, $h_\ell$ is a
nonleading coefficient of $F$, so $0<|h_\ell|<\rho_{\min}/2$ and
$\prod_pp^{2t_p}=h_\ell^2<\rho_{\min}^2\le n_j$; hence some $p$ has
$\operatorname{ord}_p(n_j)>2t_p$.  Taking the least such $p$, which depends
only on $n_j$, the $\alpha_j$ of one norm are counted at one prime.

3.  The first case is item~2 with $N=1$, and $d\ge8\rho_{\min}^2$ gives one
norm by the remark before the proposition.  If $\ell=d$, then
$v+H=A(y^d)$ with $A(z)=h_dz+v$, and \cite[Lemma~5.1]{coefficient-mass-sectors} with $t=0$
leaves at most two roots above the axis.
\end{proof}

In the Pell family after Proposition~\ref{prop:gaussheavy}, $\ell=2$, and
the three points have norms $4c^2+1$, twice, and $4c^2$: item~2 is attained,
and the bound $K\le\ell$ of item~3 fails on two norms.  Item~3 is attained
at every scale: for $w=a+ci$ on the Pell equation, $y=i\bar w=c+ai$ gives
$y^3+y=-i\,\overline{(w^3-w)}=-(8c^3+2c)$, so $(y^3+y)^2$, with nonleading
coefficients $2$ and $1$, takes the value $(8c^3+2c)^2$ at the two points
$\pm c+ai$ of one norm.  So three points of one norm need $h_1=h_2=0$, and
if $h_1\ne0$ the $\alpha_j$ have distinct norms.  The bound $K\le\ell$ is
not a constant, as $\ell$ is unbounded on one norm: if $k\ge1$ and
$a^2-3c^2=1$ with $c>\binom{2k}k$, then $(y^3+y)^{2k}=\sum_j\binom{2k}jy^{2k+2j}$, with
$\ell=2k$, is $(8c^3+2c)^{2k}$ at $\pm c+ai$, and $(y^3-y)^{2k}$, with
$\ell=2k$, is $(-1)^k(8c^3+2c)^{2k}$ at $\pm a+ci$, $2ci$; the nonleading
coefficients of both are below $c\le\rho_{\min}/2$.  In both the points of one
norm are $\alpha$ and $-\bar\alpha$, and $H$ is even.  This is forced: a
symmetry of $\mathbb Z[i]$ relating two points of a level set is a symmetry
of $H$.

\begin{proposition}[Symmetries of a level set]\label{prop:gausssym}
In the setting of Proposition~\ref{prop:gaussheavy}, item~2, write
$H=\sum_{k=1}^dh_ky^k$ with $h_d\ne0$, and call the eight numbers
$u\alpha_j$, $u\bar\alpha_j$ with $u^4=1$ the \emph{class} of $\alpha_j$.
\begin{enumerate}
\item If $\alpha\in\mathbb Z[i]$ has $|\alpha|\ge\rho_{\min}$ and
$H(\alpha)=H(u\alpha)=-v$ for a unit $u\ne1$, then $H(uy)=H(y)$: $H$ is
even if $u=-1$, and $H\in\mathbb Z[y^4]$ if $u=\pm i$.
\item A class contains at most one of the $\alpha_j$ if $H$ is not even, at
most two if $H\notin\mathbb Z[y^4]$, and at most four.
\item Let every $\alpha_j$ have norm $n$; with $t_p$, $s_p$ as in
Proposition~\ref{prop:gaussprimes}, let every prime $p\equiv1\pmod4$
dividing $n$ have $\operatorname{ord}_p(n)>2t_p$, and let $m$ of them have
$\operatorname{ord}_p(n)\ne2s_p$.  Then the $\alpha_j$ lie in at most
$\max(2^{m-1},1)$ classes.  In particular, if $n$ has at most one prime
factor $p\equiv1\pmod4$, then $K=1$ unless $H$ is even, and $K\le2$ unless
$H\in\mathbb Z[y^4]$.
\end{enumerate}
\end{proposition}

\begin{proof}
1.  Put $\rho=|\alpha|$, $x=1/\rho$ and $b=\max_{0<k<d}|h_k|$, so
$b<\rho_{\min}/2\le\rho/2$.  The polynomial
$W(y)=H(y)-H(uy)=\sum_k(1-u^k)h_ky^k$ has Gaussian-integer coefficients and
$W(\alpha)=0$.  Suppose $W\ne0$, of degree $M$.  Every $h_k$ with $k<M$ is
nonleading, so
\[
  |1-u^M|\rho^M\le|(1-u^M)h_M\alpha^M|
  =\Bigl|\sum_{k<M}(1-u^k)h_k\alpha^k\Bigr|
  \le b\sum_{k<M}|1-u^k|\rho^k.
\]
If $u=-1$, $|1-u^k|$ is $2$ for odd $k$ and $0$ for even $k$, so $M$ is
odd and the right side is at most
$2b\rho^M(x^2+x^4+\cdots)<\rho^{M+1}/(\rho^2-1)<2\rho^M$, as $\rho>2$.  If $u=\pm i$, $|1-u^k|$ is
$0,\sqrt2,2,\sqrt2$ for $k\equiv0,1,2,3\pmod4$, so $4\nmid M$, and grouping
the $k<M$ by $M-k\bmod4$ bounds the right side by
$|1-u^M|\,b\rho^M\sigma$, $\sigma=(\sqrt2x+x^2+x^4)/(1-x^4)$: this is the
case $M\equiv3$, and for $M\equiv1,2$ the weights of $x,x^2,x^3,x^4$ are
$(0,1,\sqrt2,1)$ and $(1/\sqrt2,0,1/\sqrt2,1)$, where $x^3\le x$.  Norms of
Gaussian integers above $4$ are $5$ or at least $8$.  If $\rho^2\ge8$,
$b\sigma<(\sqrt2+x+x^3)/(2(1-x^4))<0.93$; if $\rho^2=5$, then
$\rho_{\min}^2=5$, so $b\le1$ and $b\sigma<0.91$.  Either way the right side
is below $|1-u^M|\rho^M$, a contradiction.  So $W=0$.

2.  The units $u$ with $H(uy)=H(y)$ form a group $U$ of order $1$, $2$ or
$4$ as $H$ is not even, is even but not in $\mathbb Z[y^4]$, or is in
$\mathbb Z[y^4]$.  Fix $j$ and let $S$ be the set of units $u$ with
$H(u\alpha_j)=-v$.  Then $1\in S$, $U\subseteq S$, and if $u,u'\in S$,
item~1 at $u\alpha_j$, of modulus $\rho_j\ge\rho_{\min}$, gives $u'/u\in U$;
so $S=U$.  As $v+H$ has real coefficients, $u\bar\alpha_j$ is a root of it
if and only if $\bar u\alpha_j$ is, so its roots in the class of $\alpha_j$
are the $w$ and $\bar w$ with $w\in U\alpha_j$, and at most $|U|$ of them
lie above the axis.

3.  Let $p\equiv1\pmod4$ divide $n$, $e=\operatorname{ord}_p(n)$, and
$p=\pi\bar\pi$ with $\operatorname{ord}_p(\pi)>0$ under the embedding of
Proposition~\ref{prop:gaussprimes}; as $\pi$ and $\bar\pi$ are coprime,
$\operatorname{ord}_p(\bar\pi)=0$ and $\operatorname{ord}_p(\pi)=1$, so $\operatorname{ord}_p(\alpha_j)$ and $\operatorname{ord}_p(\bar\alpha_j)$
are the exponents of $\pi$ and $\bar\pi$ in $\alpha_j$.  As in the proof of
Proposition~\ref{prop:gaussprimes}, item~2, since $e>2t_p$ one of them
exceeds $t_p$ and so equals $s_p$, and they sum to $e$; so the exponent of
$\pi$ in $\alpha_j$ is $s_p$ or $e-s_p$.  The exponents of $1+i$ and of each
prime $q\equiv3\pmod4$ in $\alpha_j$ are $\operatorname{ord}_2(n)$ and
$\operatorname{ord}_q(n)/2$.  So $\alpha_j$ is determined up to a unit by a
choice between $s_p$ and $e-s_p$ at each of the $m$ primes with
$e\ne2s_p$, and $\bar\alpha_j$, in the same class, by the opposite choices.
Hence there are at most $\max(2^{m-1},1)$ classes, one if $m\le1$, and
item~2 gives the last statement.
\end{proof}

The pairs above attain item~2 with $H$ even, at every scale.  So three
points of one norm need, besides $h_1=h_2=0$, that $H\in\mathbb Z[y^4]$
(and then $\deg H\ge16$, Proposition~\ref{prop:gaussfour} below), or
that $n$ has two primes $p\equiv1\pmod4$ with
$\operatorname{ord}_p(n)\ne2s_p$ (three if $H$ is not even), or one with
$\operatorname{ord}_p(n)\le2t_p$.  Item~1 needs
the small coefficients: $H=y^5+y^4+6y^3+6y^2+25y$ is $-25$ at $\pm(1+2i)$
and is not even.

A class meeting an axis or a diagonal holds at most two points above the
axis, so by item~2 a class holds three or four of the $\alpha_j$ only if
$H=A(y^4)$ with $A\in\mathbb Z[z]$ and the class meets neither; then
$z=\alpha_j^4$ is a Gaussian integer off the real axis with $A(z)=-v$ real,
and the nonleading coefficients of $A$ lie below $|z|^{1/4}/2$.  If $H\in\mathbb
Z[y^4]$ and the $\alpha_j$ of one norm are three, one of them lies on neither
an axis nor a diagonal, since at most two Gaussian integers above the axis of one norm lie on the
axes and diagonals.

\begin{proposition}[Four points of one class]\label{prop:gaussfour}
In the setting of Proposition~\ref{prop:gaussheavy}, item~2, let
$H\in\mathbb Z[y^4]$ and let some $\alpha_j$ lie on neither axis nor
diagonal.  Then $\deg H\ge16$.
\end{proposition}

\begin{proof}
Write $H=A(y^4)$ with $A=\sum_{k=1}^Da_kz^k$ and $a_D\ne0$, so
$\deg H=4D$, and put $b=\max_{0<k<D}|a_k|$ ($b=0$ if $D=1$), $\alpha=\alpha_j=a+ic$,
$\rho=|\alpha|$ and $n=\rho^2$; then $b<\rho_{\min}/2\le\rho/2$.  Let
$\beta=\alpha^2=\mu+i\nu$ and $z=\alpha^4=\beta^2=X+iY$, so
$X=\mu^2-\nu^2$, $Y=2\mu\nu$ and $N=|z|=\mu^2+\nu^2=n^2$.  If $Y=0$ then
$(\alpha/\bar\alpha)^4=z/\bar z=1$, so $\alpha/\bar\alpha\in\{\pm1,\pm i\}$
and $\alpha$ lies on an axis or a diagonal; so $Y\ne0$.  The integers
$\mu\pm\nu=a^2\pm2ac-c^2$ are nonzero, as $c\ge1$ and $t^2\pm2t-1$ has no
rational root, and $(\mu+\nu)^2+(\mu-\nu)^2=2n^2$, so
\begin{equation}\label{eq:gaussfourX}
  |X|=|\mu+\nu|\,|\mu-\nu|\ \ge\ \max(|\mu+\nu|,|\mu-\nu|)\ \ge\ n .
\end{equation}
Suppose $D\le3$ and set $a_k=0$ for $k>D$.  As $v$ is real,
$\operatorname{Im}A(z)=0$; here $\operatorname{Im}z^2=2XY$ and
$\operatorname{Im}z^3=Y(3X^2-Y^2)=Y(4X^2-N^2)$, so dividing by $Y$,
\[
  a_1+2a_2X+a_3(4X^2-N^2)=0 .
\]
If $D=1$ this says $a_1=0$, against $a_D\ne0$.  If $D=2$ it gives $2|X|\le|a_1|\le b<\rho/2<n$,
against \eqref{eq:gaussfourX}.  Let $D=3$, put $j=2|X|-N$ and
$\sigma=X/|X|$, so $2X=\sigma(N+j)$ and
$4X^2-N^2=j(2|X|+N)=j(2N+j)$.  Then
$|a_3j|(2|X|+N)=|a_1+2a_2X|\le b(1+2|X|)$, so $|j|\le|a_3j|\le b$, and
\[
  N(2a_3j+\sigma a_2)=-(a_3j^2+\sigma a_2j+a_1),
\]
where the right side has modulus at most $2b^2+b<N$, as $b<\rho/2=N^{1/4}/2$.
So both sides vanish: $\sigma a_2=-2a_3j$ and $a_1=a_3j^2$.  If $j=0$ then
$4X^2=N^2$ and $4Y^2=3N^2$, which is impossible; so $1\le|j|\le|a_1|\le
b<\rho/2$.  But $j$ is $\mu^2-3\nu^2=n^2-4\nu^2$ if $X>0$ and
$\nu^2-3\mu^2=n^2-4\mu^2$ if $X<0$, so $j=(n-2w)(n+2w)$ with $w\ge0$ an
integer, and $j\ne0$ gives $|j|\ge n+2w\ge n>\rho/2$, a contradiction.  So
$D\ge4$.
\end{proof}

The last step is where $z$ being a fourth power, and not only a square,
enters: it makes $N=|\beta|^2$ a square.  If $\beta=a+ci$ with $a^2-3c^2=1$,
then $\beta^2-1=2c^2+2aci=2ci\bar\beta$, so $A(z)=z(z-1)^2=z^3-2z^2+z$ is
$-4c^2(a^2+c^2)^2$ at $z=\beta^2$, real, with nonleading coefficients $-2$,
$1$ below $|z|^{1/4}/2$ once $|z|=a^2+c^2>256$; here $A(y^2)=(y^3-y)^2$ and
the $\beta$ are the Pell points after Proposition~\ref{prop:gaussheavy}.  So
the sizes of the coefficients against $|z|$ do not bound $\deg A$ by
themselves.  Whether the number of Gaussian points above the axis in
one level set of a polynomial with small nonleading coefficients is bounded
independently of $\rho_{\min}$, whether one norm carries at most two of
them, and whether $\Lambda(F)\ge cK^2\log R$ holds for all integer
multiples at Gaussian roots in a wide sector, remain open.

\paragraph{Data availability.}
No datasets were generated or analyzed in this study.

\paragraph{Computer assistance.}
No proof in this paper relies on computation.  The value $-270400$ after
Proposition~\ref{prop:gaussheavy} is a direct evaluation, as are the values
$P'(0)=6$ and $P(0)$ in Proposition~\ref{prop:gausscong}, item~3, and the
expansions of $P'(0)$ in Proposition~\ref{prop:gaussline}, items~2 and~3.
Exact computations in the accompanying source check the statements and
examples; these checks are not used in any proof.

\paragraph{Use of AI tools.}
Large language model assistants (Claude, Anthropic) were used in drafting and
revising the text, the proofs and the verification code.  The author has checked the mathematics and is responsible for
the content.

\paragraph{Funding and competing interests.}
The author received no specific funding for this work and has no competing
interests to declare.

\begin{thebibliography}{99}

\bibitem{coefficient-mass-sectors}
B.~Pham,
\newblock Coefficient mass of polynomial multiples in sectors,
\newblock preprint, 2026, \url{https://github.com/bangyen/coefficient-mass}.

\bibitem{neukirch}
J.~Neukirch,
\newblock \emph{Algebraic Number Theory},
\newblock Grundlehren der mathematischen Wissenschaften 322, Springer, 1999.

\end{thebibliography}

\end{document}
